The common two-hour baseline contains 440 person-hours and has mean remaining SET excess $10.57\,^{\circ}\mathrm C$. Table~\ref{tab:deployment-results} reports the revised local-delivery calculation. At equal $68.2\%$ coverage, A36 uses the least electricity and leaves the lowest remaining discomfort among the six nominal benchmarks. Across all 24 candidates, the sampled energy--discomfort frontier consists of A3, A9, A18, A27, A36, A45 and A51. A final count depends on required coverage and residual discomfort; the finite sweep is not a global optimum.

\begin{table}[!ht]
\centering\small
\caption{Two-hour nominal results ($\beta=0.5$, $\chi=1$). Efficiencies use person$\cdot{}^{\circ}$C$\cdot$h/kWh. All six regular benchmarks have $68.2\%$ direct coverage; redistribution changes person-time coverage.}
\label{tab:deployment-results}
\begin{tabular}{@{}lrrrr@{}}
\toprule
Plan & Electricity (kWh) & Person-time coverage & Remaining excess ($^{\circ}$C) & $\eta$\\
\midrule
A36 & 18.0 & 68.2\% & 9.12 & 35.49\\
B24 & 57.6 & 68.2\% & 9.55 & 7.75\\
C18 & 90.0 & 68.2\% & 9.54 & 5.03\\
AB30 & 37.8 & 68.2\% & 9.25 & 15.39\\
AC27 & 54.0 & 68.2\% & 9.22 & 10.99\\
ABC24 & 63.9 & 68.2\% & 9.38 & 8.21\\
AB30-R & 37.8 & 76.2\% & 8.99 & 18.39\\
\bottomrule
\end{tabular}
\end{table}

\noindent\textbf{Allocation and coverage.} AB30-R relocates the same 18 A and 12 B units from equal zone counts to $(3,4,0)$, $(12,3,0)$ and $(3,5,0)$. Electricity stays at 37.8 kWh and area coverage at $68.2\%$, while person-time coverage rises to $76.2\%$. Remaining excess falls from $9.25$ to $8.99\,^{\circ}\mathrm C$ and efficiency rises from $15.39$ to $18.39$. This isolates the value of placing a fixed inventory closer to demand. Scaling all footprints by $0.75/1/1.25$ gives A36 remaining excess $9.53/9.12/8.70\,^{\circ}\mathrm C$. A36 remains the lowest-energy and lowest-residual benchmark at these common footprint multipliers; independently uncertain device footprints are not exhausted by this test.

\noindent\textbf{Why delivery assumptions matter.} Across the 15 tested combinations $\beta=0,0.25,0.5,0.75,1$ and $\chi=0,0.5,1$, A36 retains the highest efficiency among the six equal-coverage benchmarks. The lowest remaining discomfort is less stable: with $\chi=0$, AC27 becomes the lowest-residual benchmark for $\beta\ge0.5$; it also does so at $(\beta,\chi)=(1,0.5)$. At $(0.5,0)$ its remaining excess is $9.02\,^{\circ}\mathrm C$, compared with A36's $9.06\,^{\circ}\mathrm C$, but it uses 54 rather than 18 kWh. Thus local delivery can change the quality ordering without changing the efficiency ordering. The result determines a tradeoff for the specified requirements, not a universal preferred equipment type.

Changing A's assumed circulation from $0.3$ to $0.6$ to $1.2$ kg dry air/s gives A36 remaining excess $9.18/9.12/9.09\,^{\circ}\mathrm C$ and efficiency $34.06/35.49/36.22$. A36 keeps the lowest residual and highest efficiency among the six benchmarks at those settings. These checks cover a stated parameter grid; they do not establish robustness to every combination of weather, flow, coverage and passenger distribution.
